% No preamble. Prefix wcr-. Fixed reference, 17 October 2026.
\section{Direct Wilson-to-continuum heat comparison}
\label{wcr-section}

We now evaluate the fixed reference by its actual vertices and diagonal
heat kernels. No general one-loop universality theorem is imported.
Let \(h_{W,v}(s)\) and \(h_{W,g}(s)\) be the affine Cartan coefficients
of the Wilson background-Feynman vector base and the full scalar
covariant Laplacian on \(\mathbb Z^4\). Normalize
\[
 \tau(H^2)=1,\qquad C_H=\operatorname{Tr}_{\rm ad}
       (\operatorname{ad}_H^2),\qquad \alpha_g=\frac{C_H}{48\pi^2}.
\]

\begin{theorem}[Fixed-reference heat coefficient]
\label{wcr-main}
For \(s\ge1\),
\begin{equation}
 |h_{W,v}(s)-10\alpha_g|
 +|h_{W,g}(s)+\alpha_g/2|\le Cs^{-1/2}.
 \label{wcr-fixed-bound}
\end{equation}
Together with the actual-base comparison, this gives on every
applicable joint infinite structural system
\begin{equation}
 h_{M,v}(s)-2h_g(s)=11\alpha_g+O(s^{-1/2}),\qquad
 \int_1^{K^2}\bigl(h_{M,v}(s)-2h_g(s)\bigr)\frac{ds}{s}
       =22\alpha_g\log K+O(1),
 \label{wcr-middle}
\end{equation}
uniformly in the old classical depth \(M\).
\end{theorem}

\subsection{The exact Wilson density fixes the leading vertices}
\label{wcr-density-section}

The normalized unit-lattice action is
\[
 S_W(U)=\sum_{x,\mu<\nu}
 \left[1-\tfrac12\tau(P_{x,\mu,\nu}+P_{x,\mu,\nu}^{-1})\right].
\]
Each unoriented plaquette is counted once. Its left-relative Hessian and
scalar covariant gradient define
\[
 {\cal B}_W(a)=S_W''(e^{ia})+D_aD_a^T,\qquad
 L_W(a)=D_a^TD_a,\qquad
 (D_a\phi)_\mu(x)=\operatorname{Ad}_{e^{ia_\mu(x)}}\phi(x+e_\mu)-\phi(x).
\]
Transposes use the fixed bilinear counting pairing.

Introduce a physical mesh \(\epsilon\), with link angle \(\epsilon A\).
For the physical jet
\[
 z=(a,b,c,e)
 =\bigl(A_\mu,A_\nu,\nabla_{\epsilon,\mu}A_\nu,
                         \nabla_{\epsilon,\nu}A_\mu\bigr)
\]
at one plaquette, its holonomy is exactly
\begin{equation}
 P_\epsilon(z)=e^{i\epsilon a}e^{i\epsilon b+i\epsilon^2c}
              e^{-i\epsilon a-i\epsilon^2e}e^{-i\epsilon b}.
 \label{wcr-exact-plaquette}
\end{equation}
Since \(P_0=I\) and \(\partial_\epsilon P_0=0\),
\[
 F_\epsilon=\frac{P_\epsilon-I}{\epsilon^2},\qquad
 {\cal L}_\epsilon=-\tfrac12\tau(F_\epsilon^2P_\epsilon^{-1})
\]
extend jointly holomorphically to \(\epsilon=0\).
The identity \(P+P^{-1}-2I=(P-I)^2P^{-1}\) gives the exact action
\begin{equation}
 S_\epsilon(A)=\epsilon^4\sum_{x,\mu<\nu}
                 {\cal L}_\epsilon(J_\epsilon A).
 \label{wcr-exact-density}
\end{equation}
Direct differentiation of \eqref{wcr-exact-plaquette} gives
\begin{equation}
 F_0=i(c-e)-[a,b],\qquad
 {\cal L}_0=\tfrac12\tau\bigl(c-e+i[a,b]\bigr)^2.
 \label{wcr-density-limit}
\end{equation}
On each fixed bounded complex jet polydisc,
\({\cal L}_\epsilon-{\cal L}_0\), together with every fixed field
derivative, is \(O(\epsilon)\). Joint holomorphy and Cauchy's formula
supply this estimate before any momentum limit is taken.

On plane waves with bounded complex physical momenta,
\((e^{i\epsilon p_\mu}-1)/\epsilon=ip_\mu+O(\epsilon)\)
jointly analytically. Equations \eqref{wcr-exact-density} and
\eqref{wcr-density-limit} therefore identify the continuum action as
\begin{equation}
 S_c(A)=\tfrac12\int\sum_{\mu<\nu}\tau(F_{\mu\nu}(A)^2)\,dx,\qquad
 F_{\mu\nu}=\partial_\mu A_\nu-\partial_\nu A_\mu+i[A_\mu,A_\nu].
 \label{wcr-continuum-action}
\end{equation}
In particular, for the ordinary action vertices,
\begin{equation}
 \Gamma_n^W(\epsilon p_1,\ldots,\epsilon p_n)
 =\epsilon^{4-n}\Gamma_n^c(p_1,\ldots,p_n)
                  +O(\epsilon^{5-n}),\qquad n=2,3,4.
 \label{wcr-engineering}
\end{equation}
This follows by differentiating the exact density in plane-wave
amplitudes; it holds jointly with the fixed momentum derivatives used
below. In midpoint coordinates the bare quadratic correction is fourth
order. Source-site phases may give third-order bare cross entries.
The full flat Feynman symbol is exactly \(\delta(k)I\) and has the
fourth-order correction needed below, in either convention.

The physical left-relative chart is
\[
 \Xi_{\epsilon,A}(B)
   =\frac{\log(e^{i\epsilon B}e^{i\epsilon A})}{i\epsilon}
   =A+B+O(\epsilon).
\]
Its first \(B\)-derivative is \(I+O(\epsilon)\), its second is
\(O(\epsilon)\). The exact relative Hessian includes both terms
\[
 S''(A)[J_AB,J_AC]+S'(A)[T_A(B,C)].
\]
Thus it has the same continuum limit; the current term is retained.
At first background order its local corrections multiply
\(\Gamma_2\), hence have momentum order at least two.
At second background order the corrections have order at least one
from \(\Gamma_3\), or two from \(\Gamma_2\) and the current term.
They are within the remainders in \eqref{wcr-engineering}.

Likewise,
\[
 \epsilon^{-1}D_{\epsilon A}
 =\epsilon^{-1}
   \bigl(\operatorname{Ad}_{e^{i\epsilon A_\mu}}
                          \mathrm{shift}_{\epsilon e_\mu}-I\bigr)
 \longrightarrow D_{c,A,\mu}=\partial_\mu+i\operatorname{ad}_{A_\mu}
\]
jointly on these jets. Integration by parts in
\eqref{wcr-continuum-action} shows that the limiting vector and scalar
operators are
\begin{equation}
 ({\cal B}_c(A))_{\mu\nu}
 =-D_{c,A}^2\delta_{\mu\nu}-2i\operatorname{ad}_{F_{\mu\nu}},
 \qquad L_c(A)=-D_{c,A}^2.
 \label{wcr-continuum-operators}
\end{equation}
The opposite consistent adjoint sign convention gives identical heat
traces. The coefficient of the flat Laplacian is one.

\subsection{An integrable comparison of the two-point diagrams}
\label{wcr-comparison-section}

Use the same source-site Fourier convention for lattice and continuum.
For either scalar or vector their flat symbols are
\[
 B_W(k)=\delta(k)I,\qquad B_c(k)=|k|^2I,\qquad
 \delta(k)=4\sum_\mu\sin^2(k_\mu/2).
\]
On a fixed small ball with \(\rho=|k|\), their difference and its
derivatives through order two are \(O(\rho^{4-r})\).
Define \(V_a(k,p)\) and \(W_{ab}(k;p,-p)\) as the first and second
background vertices, as in the actual-base comparison.
The joint analytic expansion above gives
\begin{align}
 \|\partial_p^r(V_W-V_c)(k,p)|_{p=0}\|
 +\|\partial_p^r(V_W-V_c)(k+p,-p)|_{p=0}\|
     &\le C\rho^{2-r},\notag\\
 \|\partial_p^r(W_W-W_c)(k;p,-p)|_{p=0}\|
     &\le C\rho^{1-r},\qquad r\le2.
 \label{wcr-vertex-errors}
\end{align}
Here \(V_c\) is homogeneous of degree one and \(W_c\) is constant.
The \(r=2\) bound \(C/\rho\) on the second line is deliberately loose:
the corresponding lattice derivative is analytic and bounded.
For either baseline we can use
\(\|\partial_p^rV\|\le C\rho^{1-r}\) and
\(\|\partial_p^rW\|\le C\rho^{-r}\).

The exact total heat Hessian kernel is
\begin{align}
 {\cal H}_{ab}(s,p)
 &=-s\int\operatorname{tr}
       (e^{-sB(k)}W_{ab}(k;p,-p))\,\frac{dk}{(2\pi)^4}
       \notag\\
 &\quad+s^2\int_0^1\int\operatorname{tr}
 \left(e^{-s(1-v)B(k+p)}V_a(k,p)
       e^{-svB(k)}V_b(k+p,-p)\right)
                         \frac{dk}{(2\pi)^4}\,dv.
 \label{wcr-diagram}
\end{align}
The lattice integral is over the Brillouin zone, the continuum integral
over \(\mathbb R^4\). This identity is obtained from compact-source
Duhamel variations, or defines their absolutely convergent two-point
kernel. In the continuum Gaussian smoothing controls the first-order
differential insertions at each \(s>0\). Only total-trace cyclicity is
used.

For \(r\le2\), direct differentiation or Duhamel gives
\begin{align}
 \|\partial_k^rE(t,k)\|
   &\le C\rho^{-r}(1+t\rho^2)^r e^{-ct\rho^2},\notag\\
 \|\partial_k^r(E_W-E_c)(t,k)\|
   &\le Ct\rho^{4-r}(1+t\rho^2)^r e^{-ct\rho^2}.
 \label{wcr-heat-errors}
\end{align}
These estimates hold at \(t=0\) as well; there is no inverse heat time.

After two external derivatives, the tadpole vertex difference is bounded
by \(C\rho^{-1}e^{-cs\rho^2}\). Its integral with the prefactor \(s\)
is \(O(s^{-1/2})\) in four dimensions.
The heat-difference part is bounded by \(Cs\rho^4e^{-cs\rho^2}\)
before that prefactor, giving \(O(s^{-2})\).

Telescope the bubble product into four differences.
Replacing a vertex gives combined degree \(2+1=3\). Distributing
two external derivatives among all factors therefore gives
\[
 C\rho P(s\rho^2)e^{-cs\rho^2}.
\]
Its integral times \(s^2\) is \(O(s^{-1/2})\).
Replacing a heat factor instead gives
\[
 Cs\rho^4P(s\rho^2)e^{-cs\rho^2},
\]
whose integral times \(s^2\) is \(O(s^{-1})\).
The estimates are uniform in \(v\in[0,1]\): the two heat times sum to
\(s\), and \eqref{wcr-heat-errors} has no endpoint singularity.
Away from the low-momentum ball, both the lattice and continuum
integrals are a polynomial in \(s\) times \(e^{-cs}\).
The continuum tail uses its polynomial vertices and Gaussian decay.
Thus
\begin{equation}
 |\partial_p^2({\cal H}_{W}-{\cal H}_{c})(s,0)|
       \le Cs^{-1/2},\qquad s\ge1.
 \label{wcr-kernel-bound}
\end{equation}
Equivalently one can rescale \(k=\xi/\sqrt{s}\);
the displayed bounds provide the needed common integrable majorant.

The lattice gauge-invariant diagonal heat density and its zero linear
color jet give
\begin{equation}
 h(s)=\tfrac14\partial_{p_\nu}^2
              {\cal H}_{\alpha H,\alpha H}(s,0),\qquad\alpha\ne\nu.
 \label{wcr-density-factor}
\end{equation}
Continuum source kernels may contain contact distributions, so we do
not assume absolute continuum spatial moments or invoke a continuum
reanchoring theorem. The direct transverse diagram calculation below
identifies one quarter of the continuum derivative with its explicit
magnetic density. This and \eqref{wcr-kernel-bound} compare the affine
coefficients with error \(Cs^{-1/2}\).

\subsection{Exact continuum coefficient and sign}
\label{wcr-value-section}

Let \(q\) range over all eigenvalues of \(\operatorname{ad}_H\) on the
complexified adjoint space, with both charge signs and their
multiplicities. For one constant curvature \(F_{12}=fH\), the scalar
diagonal heat kernel in this charge sector is
\begin{equation}
 k_q(s;f)=(4\pi s)^{-2}\frac{qfs}{\sinh(qfs)}.
 \label{wcr-scalar-value}
\end{equation}
Indeed Fourier transformation in \(x_2\) gives oscillators with
eigenvalues \(|qf|(2n+1)\); integration over their centers contributes
\(|qf|/(2\pi)\). Summing the levels and multiplying by the two free
directions proves \eqref{wcr-scalar-value} as a pointwise density.

The spin matrix in \eqref{wcr-continuum-operators} has eigenvalues
\(0,0,2qf,-2qf\), so the vector density is
\[
 k_q(s;f)\{2+2\cosh(2qfs)\}.
\]
Expanding at \(f=0\) for each fixed \(s>0\) gives
\begin{equation}
 [f^2]k_q=-\frac{q^2}{96\pi^2},\qquad
 [f^2](\hbox{vector density})=\frac{10q^2}{48\pi^2}.
 \label{wcr-one-charge}
\end{equation}
This uses \(x/\sinh x=1-x^2/6+O(x^4)\) and
\(2+2\cosh(2x)=4+4x^2+O(x^4)\).
Complexification has the same trace as the original real color space,
so summing charges introduces no extra factor two.
Equations \eqref{wcr-kernel-bound}--\eqref{wcr-one-charge}
prove \eqref{wcr-fixed-bound}.

For the required direct identification of the continuum derivative,
take an external
\(a_\mu\) with momentum \(p=p_\nu e_\nu\), \(\mu\ne\nu\).
The scalar continuum first vertex is \(2qk_\mu\), while its second
vertex \(2q^2\) has no external derivative. The bubble is
\[
 2q^2s(4\pi s)^{-2}\int_0^1e^{-sv(1-v)p^2}\,dv.
\]
Its second \(p\)-derivative is \(-q^2/(24\pi^2)\), giving the scalar
value in \eqref{wcr-one-charge} after the factor \(1/4\).
The vector orbital term is four copies of this scalar term.
The two spin insertions have trace product \(8q^2p^2\), adding
\(q^2/(4\pi^2)\) to the affine coefficient and yielding the vector
value above. Mixed spin/orbital terms have zero Lorentz trace.
No spatial-moment assertion about contact distributions is needed.

For the determinant combination
\(\log\det G-\tfrac12\log\det V\), the proper-time quadratic density
is \(\tfrac12h_v-h_g\). Extraction with \(\beta=2[t^2]\) therefore
gives \(h_v-2h_g\), with this sign and no further multiplier.
Together with the actual-base bound, its middle band is
\(11\alpha_g\log(K^2)+O(1)\), proving \eqref{wcr-middle}.
For \(\tau=\operatorname{Tr}_{\rm fund}/N\) on \(SU(N)\),
\(C_H=2N^2\).

This is a positive-time density comparison. It does not assert the
existence of separate massless infinite-volume determinants at nonzero
magnetic field. Actual soft factors, their joint determinant identities
and bounded outer-band corrections must still be assembled in their
proved domains; no finite-regulator trajectory or continuum existence
is inferred from this reference computation alone.

